\chapter{Substance Use, Dependence, and Withdrawal}

\section*{Definition}
Substance withdrawal is a group of physical and psychological symptoms that can occur when a person who has developed physiological dependence abruptly stops or substantially reduces a substance. Physical dependence is an adaptation of the body and can occur with prescribed medicines; it is not the same as addiction or a substance use disorder, which involves problematic or compulsive use despite harmful consequences.

\section*{When to See a Doctor}

\begin{itemize}
 \item Call emergency services for a seizure, severe confusion or hallucinations, delirium, difficulty breathing, chest pain, collapse, severe agitation, suicidal thoughts, or suspected overdose.
 \item Seek urgent medical assessment for withdrawal after heavy or prolonged alcohol, benzodiazepine, or barbiturate use; repeated vomiting or diarrhea with dehydration; severe depression; pregnancy; or withdrawal in a person with serious medical or psychiatric illness.
 \item Do not attempt to stop alcohol or benzodiazepines abruptly without medical advice, even if previous withdrawal was mild.
\end{itemize}

\section*{How It Develops}
Repeated exposure to a substance can produce neuroadaptation in several brain systems. Abrupt cessation or a large dose reduction removes the substance's effects before the body has readjusted, producing symptoms that vary by substance, dose, duration, formulation, co-use of other substances, and individual health. Withdrawal may involve increased nervous-system activity, reduced mood or energy, gastrointestinal symptoms, pain, or other substance-specific effects.

\section*{Causes}
\begin{itemize}
 \item Alcohol withdrawal, which can progress to seizures or delirium tremens and can be fatal without treatment
 \item Opioid withdrawal, which commonly causes aches, abdominal cramps, diarrhea, nausea, sweating, anxiety, restlessness, and cravings; loss of tolerance after stopping increases overdose risk if opioid use is resumed
 \item Benzodiazepine withdrawal, which may be prolonged and can cause seizures, delirium, or severe anxiety; abrupt discontinuation can be dangerous
 \item Nicotine withdrawal (cigarettes, vaping)
 \item Stimulant withdrawal (amphetamine, cocaine — fatigue, depression)
 \item Cannabis withdrawal (irritability, insomnia, anxiety)
 \item Discontinuation symptoms after some antidepressants or antipsychotics; these are not always the same as a substance withdrawal syndrome and may include dizziness, sleep disturbance, nausea, movement symptoms, or return of the underlying illness
 \item Caffeine withdrawal (headache, fatigue, irritability)
 \item Antipsychotic withdrawal
 \item Barbiturate withdrawal, which can be severe and life-threatening
\end{itemize}

\section*{Related Symptoms}
\begin{itemize}
 \item Anxiety, irritability, agitation, restlessness, or craving
 \item Insomnia, nightmares, tremor, sweating, chills, or palpitations
 \item Nausea, vomiting, abdominal cramps, diarrhea, reduced appetite, or muscle aches
 \item Headache, poor concentration, fatigue, low mood, or suicidal thoughts
 \item Hallucinations, seizures, confusion, severe tremor, or delirium in complicated withdrawal
\end{itemize}

\section*{Medical Evaluation}
\begin{itemize}
 \item A clinician will assess the substance, dose, route, formulation, duration and pattern of use, timing of the last dose, previous withdrawal, prescribed medicines, alcohol and polysubstance use, medical and psychiatric history, pregnancy possibility, and suicide or overdose risk.
 \item The examination includes vital signs, mental status, hydration, neurologic findings, tremor, and signs of intoxication, infection, trauma, or coexisting illness. CIWA-Ar may support assessment of alcohol withdrawal and COWS may support assessment of opioid withdrawal, but neither score replaces clinical judgment and neither is appropriate for every substance or patient.
 \item Tests are selected according to the presentation and may include glucose, electrolytes, kidney and liver tests, blood count, pregnancy testing, ECG, toxicology testing, or tests for infection and other causes of the symptoms. A urine toxicology screen has limitations and does not by itself diagnose addiction, identify every substance, or determine withdrawal severity.
 \item Moderate-to-severe or medically complicated withdrawal requires medically supervised withdrawal management, with referral to emergency, inpatient, addiction, or specialist services according to risk. Withdrawal management should be linked to ongoing treatment for any substance use disorder.
\end{itemize}

\section*{Treatment Options}
\begin{itemize}
 \item Alcohol withdrawal is treated under medical supervision, commonly with a benzodiazepine-based regimen, supportive care, and thiamine when indicated; severe withdrawal may require hospital or intensive care.
 \item Benzodiazepine withdrawal generally requires an individualized, gradual taper supervised by the prescriber or an addiction specialist; abrupt cessation should be avoided.
 \item Opioid withdrawal symptoms may be treated with medicines such as buprenorphine, methadone, or lofexidine when clinically appropriate. For opioid use disorder, ongoing medication treatment with buprenorphine, methadone, or naltrexone is more than short-term detoxification and reduces relapse and overdose risk.
 \item Symptom-triggered or scheduled treatment is selected by trained clinicians using the substance, severity, comorbidities, and monitoring setting; scores are aids rather than automatic dosing instructions.
 \item Long-term recovery may include medication, behavioral therapy, peer or community support, harm-reduction services, and treatment of co-occurring psychiatric or medical conditions. Naloxone should be offered or prescribed when opioid overdose risk is present.
\end{itemize}

\section*{Self-Care and Home Management}
\begin{itemize}
 \item Do not try to manage alcohol, benzodiazepine, or barbiturate withdrawal alone. Contact a healthcare provider or addiction service before stopping or reducing these substances.
 \item While awaiting advice for mild symptoms, take small sips of fluid if you can keep them down and eat as tolerated; seek urgent care for repeated vomiting, dehydration, confusion, hallucinations, or worsening symptoms.
 \item Ask a clinician or pharmacist before taking sedatives, sleep medicines, opioids, or supplements during withdrawal, and do not combine opioids with alcohol or benzodiazepines.
 \item Arrange support from a trusted person and consider peer-support or community services, but do not use peer support as a substitute for emergency or medical care.
 \item If opioid use is stopped or reduced, keep naloxone available where possible and do not resume a previous dose; tolerance can fall quickly and the old dose may cause a fatal overdose.
\end{itemize}

\section*{References}
\begin{thebibliography}{99}
\bibitem{addiction_withdrawal:ref1} American Society of Addiction Medicine. The ASAM Clinical Practice Guideline on Alcohol Withdrawal Management. ASAM, 2020.
\bibitem{addiction_withdrawal:ref2} American Society of Addiction Medicine. The ASAM National Practice Guideline for the Treatment of Opioid Use Disorder: 2020 Focused Update.
\bibitem{addiction_withdrawal:ref3} Substance Abuse and Mental Health Services Administration. TIP 63: Medications for Opioid Use Disorder. Updated 2021.
\bibitem{addiction_withdrawal:ref4} National Institute on Drug Abuse. Opioids; Principles of Drug Addiction Treatment. National Institutes of Health, current online resources.
\bibitem{addiction_withdrawal:ref5} World Health Organization. Management of Alcohol Withdrawal. Current online guidance.
\bibitem{addiction_withdrawal:ref6} American Society of Addiction Medicine and partner organizations. Joint Clinical Practice Guideline on Benzodiazepine Tapering. 2025.
\end{thebibliography}
